% ---------------------------------------------------------------------------
% Modelagem — Frente 2
% Projeto SMAC03 (Grafos / UNIFEI) — Cenario 1: Analise do Mercado de FIIs
%
% ONDE ENTRA: dentro da secao 3 (Desenvolvimento), logo depois de
% "3.1 Coleta de Dados". Esta escrito como \subsection, virando "3.2 Modelagem"
% e substituindo o "3.2 Exemplo de Algoritmo" que veio do template.
%
% PREAMBULO: o documento precisa de \usepackage{amsmath}, \usepackage{amssymb}
% e \usepackage{booktabs}.
%
% Os numeros conferem com out/frente2/ e com docs/frente2-thresholds.md.
% Reproduzir com: python src/frente2.py
% ---------------------------------------------------------------------------

\subsection{Modelagem}

A modelagem proposta representa o conjunto de fundos analisados como três
grafos não direcionados e valorados, construídos sobre o mesmo conjunto de
vértices. Formalmente, cada camada constitui um grafo
$G_\ell = (V, E_\ell, w_\ell)$, com $\ell \in \{A, B, C\}$, em que $V$ é o
conjunto dos 46 Fundos de Investimento Imobiliário remanescentes após o
tratamento descrito na Seção~3.1, $E_\ell$ é o conjunto de arestas daquela
camada e $w_\ell : E_\ell \rightarrow \mathbb{R}_{+}$ associa a cada aresta a
intensidade da relação que ela representa.

A opção por três camadas, e não por um único índice agregado, decorre do fato
de que a noção de semelhança entre fundos admite mais de uma definição. Dois
fundos podem compartilhar características declaradas, podem oscilar de maneira
conjunta no mercado ou podem apresentar perfis financeiros próximos sem que
nenhuma das condições anteriores se verifique. Desse modo, tratar essas noções
como grafos distintos permite observar não apenas onde elas concordam, mas
sobretudo onde divergem, que é precisamente o ponto em que a diversificação
aparente de uma carteira se separa da diversificação efetiva.

\subsubsection{Camada A --- similaridade estrutural}

A primeira camada parte das características declaradas de cada fundo, a saber,
o segmento de atuação, a gestora responsável e os principais locatários. A
amostra é representada por uma matriz bipartida
$M \in \{0,1\}^{n \times m}$, na qual $n = 46$ corresponde ao número de fundos
e $m$ ao número de itens observados, sendo cada item um par atributo--valor,
como \texttt{segmento=Logística} ou \texttt{inquilino=Amazon}. Tem-se
$M_{iq} = 1$ quando o fundo $i$ apresenta o item $q$. A projeção dessa matriz
sobre o conjunto de fundos, dada por $M M^{\top}$, produz uma matriz cujo
elemento $(i,j)$ contabiliza os itens compartilhados pelos fundos $i$ e $j$.

A contagem simples, entretanto, atribui a todos os itens o mesmo valor
informativo, o que não se sustenta: compartilhar um segmento presente em
quinze dos quarenta e seis fundos é consideravelmente menos informativo do que
compartilhar uma gestora presente em dois. Dessa forma, adotou-se uma
ponderação por raridade, análoga ao fator de frequência inversa empregado em
recuperação de informação. Sendo $n_q$ o número de fundos que apresentam o
item $q$, define-se

\begin{equation}
  \mathrm{idf}(q) = \ln\!\left(\frac{n}{n_q}\right).
\end{equation}

A ponderação por raridade, isoladamente, ainda apresenta duas deficiências,
identificadas durante a calibragem. Em primeiro lugar, a soma do
$\mathrm{idf}$ dos itens compartilhados favorece sistematicamente os fundos
que possuem mais atributos preenchidos: o peso médio de uma aresta interna ao
segmento de logística era de $2{,}57$, contra $1{,}31$ no segmento de papel,
diferença que decorre da existência de locatários em um caso e não no outro.
Em segundo lugar, a soma depende exclusivamente de quais atributos coincidem,
e não da extensão em que os perfis se assemelham, de modo que o grafo completo
apresentava apenas vinte e dois valores distintos de peso para duzentas e
trinta e três arestas, das quais cinquenta e oito por cento eram mutuamente
empatadas por compartilharem somente o segmento.

Por conseguinte, adotou-se o índice de Jaccard ponderado pelo
$\mathrm{idf}$, que mede a fração do perfil conjunto de dois fundos que lhes
é comum. Sendo $A$ e $B$ os conjuntos de itens dos fundos $i$ e $j$,

\begin{equation}
  w_A(i,j) =
  \frac{\sum_{q \in A \cap B} \mathrm{idf}(q)}
       {\sum_{q \in A \cup B} \mathrm{idf}(q)},
\end{equation}

com $w_A(i,j) \in (0,1]$. O denominador introduz dependência também naquilo
que distingue os fundos, elevando o número de valores distintos de peso de
vinte e dois para cento e quatorze e corrigindo o desequilíbrio entre
segmentos.

Cumpre observar que o atributo referente aos locatários é aplicado apenas aos
fundos de tijolo. Fundos de papel, fundos de fundos e fundos do agronegócio
não possuem locatários, e o que a base registra nesses casos é uma descrição
genérica do lastro, idêntica para todos os fundos de um mesmo segmento. Sua
utilização como atributo faria com que os fundos de papel compartilhassem onze
itens simultaneamente, originando um subgrafo completo de peso elevado que
apenas reiteraria a informação já contida no segmento.

Além do mais, a projeção bipartida tende a produzir grafos densos, visto que
um atributo categórico compartilhado por $n_q$ fundos origina um subgrafo
completo de $n_q$ vértices. A redução por limiar de peso mostrou-se inadequada
neste caso específico: como o peso de uma aresta originada exclusivamente do
segmento compartilhado é dominado por $\ln(n/n_q)$, qualquer limiar elimina
primeiramente os segmentos mais numerosos, que são justamente aqueles de maior
interesse para a análise de concentração de risco.

Adotou-se, em substituição, uma extração de \textit{backbone} por vizinhança,
na qual cada fundo retém as $k$ arestas de maior peso que nele incidem, bem
como aquelas que com elas empatam, sendo o grafo final a união dessas
seleções. Outrossim, o critério passa a ser interno a cada vértice e,
portanto, invariante ao tamanho do grupo a que ele pertence; a manutenção dos
empates assegura que o resultado não dependa da ordem de leitura dos dados; e
nenhum fundo permanece isolado.

\subsubsection{Camada B --- co-movimento de mercado}

A segunda camada não utiliza atributos declarados, mas as séries históricas de
preços. A partir das cotações diárias de fechamento ajustado $P_i(t)$,
calcula-se o retorno logarítmico

\begin{equation}
  r_i(t) = \ln P_i(t) - \ln P_i(t-1),
\end{equation}

e, sobre esses retornos, o coeficiente de correlação linear de Pearson
$\rho_{ij}$ entre cada par de fundos. A correlação é então convertida na
métrica proposta por Mantegna (1999),

\begin{equation}
  d_{ij} = \sqrt{2\,(1 - \rho_{ij})},
\end{equation}

que satisfaz os axiomas de distância e é monotonicamente decrescente em
$\rho_{ij}$, de maneira que pares fortemente correlacionados tornam-se
próximos no espaço métrico resultante.

Sobre essa matriz de distâncias constrói-se a árvore geradora mínima, que
conecta os quarenta e seis fundos ao menor custo total e evidencia, para cada
fundo, o seu vizinho mais próximo em termos de co-movimento. A árvore,
contudo, é restrita a $n-1$ arestas e descarta ligações fortes que não lhe
pertencem. Dessa forma, o grafo da camada foi definido como a união da árvore
geradora mínima com o grafo por limiar:

\begin{equation}
  E_B = E_{\mathrm{MST}} \;\cup\; \{(i,j) : \rho_{ij} \geq \rho^{*}\}.
\end{equation}

A construção por união justifica-se pelo fato de a distribuição de correlação
da amostra ser estreita, apresentando mediana de $0{,}185$ e máximo de
$0{,}484$, de modo que um limiar isolado ou retém a quase totalidade dos pares
ou fragmenta o grafo. Com $\rho^{*} = 0{,}32$, o grafo por limiar isolado
deixaria dezoito dos quarenta e seis fundos sem aresta alguma, inviabilizando
o cálculo de centralidade de intermediação, caminho médio, diâmetro e dos
ensaios de robustez. A união preserva a conectividade e admite duas leituras
distintas para as arestas: aquelas provenientes do limiar indicam pares
efetivamente co-movidos, ao passo que aquelas provenientes apenas da árvore
indicam o vizinho mais próximo de um fundo, ainda que a correlação seja
moderada. Tal distinção é registrada como atributo da aresta e representada
graficamente por traço descontínuo.

\subsubsection{Camada C --- similaridade de perfil fundamentalista}

A terceira camada descreve cada fundo por um vetor de cinco indicadores
quantitativos, a saber: \textit{dividend yield} acumulado em doze meses, razão
entre preço e valor patrimonial, patrimônio líquido, liquidez diária e número
de ativos. O patrimônio líquido e a liquidez são transformados por
$\log_{10}$, por variarem entre milhões e bilhões, e a totalidade dos
indicadores é padronizada por escore-padrão, a fim de que nenhuma variável
domine a distância em razão de sua unidade de medida. Valores inaplicáveis,
como o número de ativos em fundos de papel e em fundos de fundos, são tratados
como ausentes e imputados pela média da respectiva variável, equivalente a
zero após a padronização.

A dissimilaridade entre dois fundos corresponde à distância euclidiana
$d_{ij}$ entre seus vetores padronizados, convertida em similaridade por

\begin{equation}
  w_C(i,j) = \frac{1}{1 + d_{ij}},
\end{equation}

que é limitada ao intervalo $(0,1]$ e decresce suavemente com a distância.
Visto que tal definição originaria um grafo completo, aplica-se a mesma
extração por vizinhança empregada na Camada A, de modo que cada fundo retém os
seus $k$ vizinhos mais similares.

\subsubsection{Calibragem dos limiares}

Os parâmetros das três camadas foram submetidos a uma análise de
sensibilidade, na qual se observou a resposta do número de arestas, da
densidade, do grau médio, do número de componentes conexos e da modularidade a
cada valor testado. Somam-se a eles os dois critérios de qualidade de dados
descritos na Seção~3.1, igualmente tratados como parâmetros do modelo.

O limiar da Camada B foi fixado no percentil 95 da distribuição de correlação
da própria amostra, igual a $0{,}315$ e arredondado para $0{,}32$, de modo que
as arestas de limiar correspondam, por definição, aos cinco por cento de pares
mais correlacionados. Para as Camadas A e C adotou-se $k = 3$, o menor valor
que mantém ambas com grau médio superior a três, compatível com o da Camada B;
na Camada A, adicionalmente, a modularidade é máxima até $k = 3$ e decresce a
partir de $k = 4$. A adoção de um mesmo critério para as duas camadas resulta
em grafos de densidade comparável, condição necessária para que a comparação
entre camadas reflita diferenças de estrutura e não de calibragem.

\begin{table}[htbp]
\centering
\caption{Limiares adotados.}
\label{tab:limiares}
\begin{tabular}{llll}
\toprule
\textbf{Limiar} & \textbf{Valor} & \textbf{Camada} & \textbf{Critério} \\
\midrule
Cobertura mínima         & 0,90 & dados & cotação efetiva em 90\% do período \\
Máximo de retornos nulos & 0,30 & dados & acima disso, preço travado \\
$k$ do \textit{backbone} & 3    & A     & grau médio acima de três \\
$\rho^{*}$               & 0,32 & B     & percentil 95 da correlação \\
$k$ da vizinhança        & 3    & C     & mesmo critério da Camada A \\
\bottomrule
\end{tabular}
\end{table}

\begin{table}[htbp]
\centering
\caption{Grafos resultantes após a calibragem.}
\label{tab:camadas}
\begin{tabular}{lrrrrrr}
\toprule
\textbf{Camada} & \textbf{Arestas} & \textbf{Densidade} & \textbf{Grau médio}
& \textbf{Comp.} & \textbf{Comun.} & \textbf{Modularidade} \\
\midrule
A & 122 & 0,118 & 5,3 & 1 & 6 & 0,574 \\
B &  68 & 0,066 & 3,0 & 1 & 5 & 0,422 \\
C &  95 & 0,092 & 4,1 & 1 & 6 & 0,664 \\
\bottomrule
\end{tabular}
\end{table}

Como validação da modelagem, comparou-se a partição obtida pelo algoritmo de
Louvain em cada camada com a classificação setorial dos fundos, por meio do
índice Rand ajustado. A Camada A apresenta concordância de $0{,}519$,
resultado esperado, dado que o segmento constitui um de seus atributos de
entrada. A Camada B apresenta $0{,}210$ e a Camada C, $0{,}122$. Nesse viés, a
divergência entre esses valores sustenta a decisão de manter as três camadas:
o co-movimento de mercado acompanha apenas parcialmente a divisão setorial, ao
passo que o perfil financeiro mostra-se praticamente ortogonal a ela. Caso
alguma das camadas fosse redundante em relação às demais, tal redundância se
manifestaria nesse indicador.

\subsubsection{Limitações da modelagem}

A modelagem está sujeita a cinco limitações que condicionam a leitura dos
resultados. Em primeiro lugar, os limiares constituem decisões de modelagem e
não parâmetros estimados, razão pela qual a análise de sensibilidade documenta
o comportamento das métricas sob valores alternativos. Em segundo lugar, o
índice de Jaccard divide pela união dos perfis, de maneira que fundos sobre os
quais se dispõe de mais informação possuem mais oportunidade de divergir: o
peso médio interno ao segmento de logística é de $0{,}158$, contra $0{,}324$
no de papel, razão pela qual a comparação de pesos entre segmentos distintos
deve ser feita com reserva, ainda que a comparação interna a um mesmo segmento
permaneça válida. Em terceiro lugar, a correlação é sensível à janela adotada
e a mudanças de regime, de modo que os resultados se referem especificamente
ao biênio 2023--2024. Em quarto lugar, a presença de aresta na Camada B não
implica correlação elevada, dado que a árvore geradora impõe $n-1$ arestas
independentemente de sua intensidade. Por fim, o atributo referente aos
locatários contribui em apenas dezenove dos quarenta e seis fundos, por
aplicar-se somente aos segmentos de tijolo.
